\documentclass[aspectratio=169,10pt]{beamer}

% Shared Beamer setup for Fall 2026 course slides.
\mode<presentation>{
  \usetheme{Madrid}
  \usecolortheme{default}
}

\definecolor{CalPolyGreen}{HTML}{154734}
\definecolor{CalPolyGold}{HTML}{BD8B13}
\definecolor{DeckInk}{HTML}{1F2933}
\definecolor{DeckMuted}{HTML}{5F6B7A}

\setbeamercolor{structure}{fg=CalPolyGreen}
\setbeamercolor{palette primary}{bg=CalPolyGreen,fg=white}
\setbeamercolor{palette secondary}{bg=DeckInk,fg=white}
\setbeamercolor{palette tertiary}{bg=CalPolyGold,fg=DeckInk}
\setbeamercolor{frametitle}{bg=CalPolyGreen,fg=white}
\setbeamercolor{title}{fg=CalPolyGreen}
\setbeamercolor{normal text}{fg=DeckInk,bg=white}
\setbeamercolor{block title}{bg=CalPolyGreen,fg=white}
\setbeamercolor{block body}{bg=black!3,fg=DeckInk}

\setbeamertemplate{navigation symbols}{}
\setbeamertemplate{footline}[frame number]
\setbeamertemplate{itemize item}{\color{CalPolyGold}$\blacktriangleright$}
\setbeamertemplate{itemize subitem}{\color{CalPolyGreen}$\bullet$}

\newcommand{\CourseNumber}{Course}
\newcommand{\CourseTitle}{Course Title}
\newcommand{\LectureNumber}{0}
\newcommand{\LectureTitle}{Lecture Title}
\newcommand{\LectureDate}{Date}
\newcommand{\CourseUnit}{Unit}

\newcommand{\coursenumber}[1]{\renewcommand{\CourseNumber}{#1}}
\newcommand{\coursetitle}[1]{\renewcommand{\CourseTitle}{#1}}
\newcommand{\lecturenumber}[1]{\renewcommand{\LectureNumber}{#1}}
\newcommand{\lecturetitle}[1]{\renewcommand{\LectureTitle}{#1}}
\newcommand{\lecturedate}[1]{\renewcommand{\LectureDate}{#1}}
\newcommand{\courseunit}[1]{\renewcommand{\CourseUnit}{#1}}

\author{Kirk Duran}
\institute{California Polytechnic State University, San Luis Obispo}
\date{\LectureDate}

\newcommand{\makedecktitle}{
  \begin{frame}[plain]
    \vspace{0.25in}
    {\usebeamerfont{title}\color{CalPolyGreen}\LectureTitle\par}
    \vspace{0.18in}
    {\large \CourseNumber{}: \CourseTitle\par}
    \vspace{0.08in}
    {\large Lecture \LectureNumber{} -- \LectureDate\par}
    \vfill
    {\small Kirk Duran \textbar{} Cal Poly, San Luis Obispo\par}
  \end{frame}
}

\newcommand{\placeholdernote}{
  \begin{block}{Placeholder}
    Replace these bullets with the final lecture material, examples, and in-class activities.
  \end{block}
}

\AtBeginSection[]{
  \begin{frame}{Roadmap}
    \tableofcontents[currentsection]
  \end{frame}
}


\graphicspath{{../assets/lecture-04/}}

% If an image file is not yet available, skip the visual slot.
\newcommand{\lectureimage}[2][1.75in]{%
  \IfFileExists{../assets/lecture-04/#2}{%
    \includegraphics[width=\linewidth,height=#1,keepaspectratio]{#2}%
  }{%
  }%
}

\setbeamersize{text margin left=0.42in,text margin right=0.42in}

\coursenumber{CS 4880}
\coursetitle{Artificial Intelligence}
\lecturenumber{4}
\lecturetitle{Uninformed Search I}
\lecturedate{August 31, 2026}
\courseunit{Search and Problem Solving}

\begin{document}

\makedecktitle

\begin{frame}[shrink=20]{Today}
  \begin{itemize}
    \item Why search is a legitimate computational model of intelligent problem solving.
    \item How to formulate a problem using states, actions, transitions, and goals.
    \item How search trees, frontiers, explored sets, and parent pointers organize exploration.
    \item How Breadth-First Search, Depth-First Search, and Iterative Deepening differ.
    \item How autonomous Micromouse robots turn an unknown maze into a search problem.
  \end{itemize}

  \begin{keyidea}
    Uninformed search asks how an agent can reason systematically about possible futures without any estimate of which unexplored state is most promising.
  \end{keyidea}
\end{frame}

\sectionframe{Part 1: Search as Intelligent Problem Solving}

\begin{frame}[shrink=20]{When Reflex Behavior Is Not Enough}
  \begin{columns}[T,onlytextwidth]
    \begin{column}{0.54\textwidth}
      \begin{block}{Reflex agent}
        percept $\longrightarrow$ immediate action
      \end{block}
      \begin{itemize}
        \item Many tasks require a sequence of actions whose consequences extend beyond the current percept.
        \item The correct first action may only become apparent after considering what can happen later.
      \end{itemize}

      \begin{block}{Problem-solving agent}
        current state $\longrightarrow$ possible futures $\longrightarrow$ plan
      \end{block}
    \end{column}

    \begin{column}{0.40\textwidth}
      % Search direction: branching future-state diagram from one current state to several successors and a goal.
      \lectureimage[2.38in]{slide4.png}
    \end{column}
  \end{columns}

  \begin{keyidea}
    Search allows an agent to reason about states it has not yet physically encountered.
  \end{keyidea}
\end{frame}

\begin{frame}[shrink=20]{Problem Abstraction Removes Irrelevant Detail}
  \begin{cpdefinition}
    A \textbf{problem formulation} is an abstract mathematical model of a task that retains the information needed to find a solution.
  \end{cpdefinition}

  \begin{columns}[T,onlytextwidth]
    \begin{column}{0.48\textwidth}
      \begin{block}{Physical navigation problem}
        Traffic, weather, buildings, passengers, road geometry, fuel, and many other details may exist in the real environment.
      \end{block}
    \end{column}

    \begin{column}{0.48\textwidth}
      \begin{block}{Search abstraction}
        \begin{itemize}
          \item State $\rightarrow$ current location
          \item Action $\rightarrow$ move to an adjacent location
          \item Goal $\rightarrow$ destination
        \end{itemize}
      \end{block}
    \end{column}
  \end{columns}

  \begin{keyidea}
    The search algorithm operates on the abstraction, not on every physical detail of the world.
  \end{keyidea}
\end{frame}

\begin{frame}[shrink=20]{Problem-Solving Agents Treat States as Atomic}
  \begin{itemize}
    \item For uninformed search, a state can be treated as a black box with a unique identity.
    \item The algorithm needs to know which actions are available, which successor states result, and whether a goal has been reached.
    \item It does not require semantic understanding of the internal meaning of a state.
  \end{itemize}

  \begin{center}
    \renewcommand{\arraystretch}{1.25}
    \begin{tabular}{l|l}
      \textbf{Domain} & \textbf{Possible meaning of a state} \\\hline
      Navigation & current location \\
      Sliding puzzle & tile configuration \\
      Theorem proving & partially constructed proof \\
      Robotics & current discrete configuration
    \end{tabular}
  \end{center}

  \begin{keyidea}
    The same search algorithm can operate across domains because the representation supplies the meaning.
  \end{keyidea}
\end{frame}

\begin{frame}[shrink=20]{Newell and Simon: Reasoning as Search}
  \begin{columns}[T,onlytextwidth]
    \begin{column}{0.55\textwidth}
      \begin{itemize}
        \item Newell and Simon's early AI work treated problem solving as search through a symbolic \textbf{problem space}.
        \item Their Logic Theorist represented theorem proving using states, operators, and goals.
      \end{itemize}

      \begin{block}{Logic Theorist, 1956}
        \begin{itemize}
          \item States: partially constructed proofs
          \item Actions: legal inference rules
          \item Goal: derive the target theorem
        \end{itemize}
      \end{block}

      \begin{keyidea}
        The claim was substantive: important forms of reasoning might be implemented as systematic search.
      \end{keyidea}
    \end{column}

    \begin{column}{0.39\textwidth}
      % Search direction: Newell and Simon / Logic Theorist historical image or theorem-proof search tree.
      \lectureimage[2.50in]{slide7.jpg}
    \end{column}
  \end{columns}

  \note{[Sources] Allen Newell and Herbert A. Simon, early work on Logic Theorist and symbolic problem solving; see also Russell and Norvig, Artificial Intelligence: A Modern Approach, search chapter.}
\end{frame}

\sectionframe{Part 2: Formalizing an Uninformed Search Problem}

\begin{frame}[shrink=20]{Our Search Model Today}
  \begin{columns}[T,onlytextwidth]
    \begin{column}{0.55\textwidth}
      \begin{description}
        \item[State space] The states relevant to the problem.
        \item[Initial state] Where the agent begins.
        \item[Goal test] Determines whether a state satisfies the objective.
        \item[Actions] $\mathrm{ACTIONS}(s)$ returns legal actions in state $s$.
        \item[Transition model] $\mathrm{RESULT}(s,a)$ returns the state produced by action $a$.
      \end{description}
    \end{column}

    \begin{column}{0.39\textwidth}
      \begin{block}{Uniform-step assumption}
        \begin{itemize}
          \item Every legal action counts as one step.
          \item We ignore variable action costs.
          \item A shortest solution has the fewest actions.
        \end{itemize}
      \end{block}

      \begin{alertblock}{Not today}
        Uniform-Cost Search, Dijkstra's algorithm, and weighted path costs.
      \end{alertblock}
    \end{column}
  \end{columns}
\end{frame}

\begin{frame}[shrink=20]{A Solution Is a Sequence of Actions}
  \begin{columns}[T,onlytextwidth]
    \begin{column}{0.54\textwidth}
      \begin{cpdefinition}
        A \textbf{solution} is an action sequence that transforms the initial state into a goal state.
      \end{cpdefinition}

      \begin{block}{Example}
        \centering
        $s_0 \xrightarrow{a_1} s_1 \xrightarrow{a_2} s_2 \xrightarrow{a_3} s_g$
      \end{block}

      \begin{itemize}
        \item Solution: $[a_1,a_2,a_3]$
        \item Solution depth: $3$
      \end{itemize}
    \end{column}

    \begin{column}{0.40\textwidth}
      % Search direction: sequence of states and actions from initial state to goal.
      \lectureimage[2.20in]{slide10.png}
    \end{column}
  \end{columns}

  \begin{keyidea}
    In today's model, a shortest solution is a goal path of minimum depth.
  \end{keyidea}
\end{frame}

\begin{frame}[shrink=20]{State-Space Graph vs. Search Tree}
  \begin{columns}[T,onlytextwidth]
    \begin{column}{0.47\textwidth}
      \begin{block}{State-space graph}
        Represents the underlying problem domain.
        \begin{itemize}
          \item vertices $=$ states;
          \item edges $=$ legal transitions.
        \end{itemize}
      \end{block}
    \end{column}

    \begin{column}{0.47\textwidth}
      \begin{block}{Search tree}
        Generated dynamically by the algorithm.
        \begin{itemize}
          \item root $=$ initial state;
          \item children $=$ generated successors.
        \end{itemize}
      \end{block}
    \end{column}
  \end{columns}

  \begin{itemize}
    \item Multiple search-tree nodes can correspond to the same state-space vertex.
    \item The tree records paths being considered, not a second physical world.
  \end{itemize}

  \begin{keyidea}
    The problem defines the graph; the algorithm generates only the portion of the search tree it needs.
  \end{keyidea}
\end{frame}

\begin{frame}[shrink=20]{Why Generate States On Demand?}
  \begin{columns}[T,onlytextwidth]
    \begin{column}{0.52\textwidth}
      \begin{itemize}
        \item Many search spaces are too large to enumerate in advance.
        \item 8-puzzle permutations: $9! = 362{,}880$.
        \item 15-puzzle permutations: $16! \approx 2.09 \times 10^{13}$.
        \item More realistic AI problems can have vastly larger or effectively unbounded spaces.
      \end{itemize}

      \begin{keyidea}
        Generate successors only when a node is expanded.
      \end{keyidea}
    \end{column}

    \begin{column}{0.42\textwidth}
      % Search direction: small generated search tree beside a huge faded implicit state space.
      \lectureimage[2.40in]{slide12.png}
    \end{column}
  \end{columns}
\end{frame}

\begin{frame}[shrink=20]{Frontier and Explored Set}
  \begin{columns}[T,onlytextwidth]
    \begin{column}{0.47\textwidth}
      \begin{block}{Frontier}
        Nodes discovered but not yet expanded.
      \end{block}
      \begin{itemize}
        \item boundary between explored and unexplored possibilities;
        \item strategy determines which frontier node is expanded next.
      \end{itemize}
    \end{column}

    \begin{column}{0.47\textwidth}
      \begin{block}{Explored set}
        States that have already been expanded.
      \end{block}
      \begin{itemize}
        \item avoids unnecessary repeated exploration;
        \item especially important when the graph contains cycles.
      \end{itemize}
    \end{column}
  \end{columns}

  \begin{keyidea}
    BFS and DFS use the same basic search machinery; they differ primarily in how the frontier is organized.
  \end{keyidea}
\end{frame}

\begin{frame}[shrink=20]{Parent Pointers Turn a Goal into a Plan}
  \begin{columns}[T,onlytextwidth]
    \begin{column}{0.54\textwidth}
      \begin{itemize}
        \item Each newly generated search node can store the node that generated it.
        \item When a goal is found, follow parent pointers backward to the initial node.
        \item Reverse the sequence to recover the forward solution path.
      \end{itemize}

      \begin{block}{Search-node information}
        \texttt{state} \qquad \texttt{parent} \qquad \texttt{action}
      \end{block}
    \end{column}

    \begin{column}{0.40\textwidth}
      % Search direction: goal node with arrows following parent pointers back to start.
      \lectureimage[2.35in]{slide14.png}
    \end{column}
  \end{columns}

  \begin{keyidea}
    Search discovers a goal; parent information reconstructs the action sequence that reaches it.
  \end{keyidea}
\end{frame}

\begin{frame}[shrink=20]{How Should We Evaluate a Search Strategy?}
  \begin{description}
    \item[Completeness] Is a solution guaranteed to be found when one exists?
    \item[Shortest-solution guarantee] Is the returned solution guaranteed to have minimum depth?
    \item[Time complexity] How many search nodes may be generated or expanded?
    \item[Space complexity] How many search nodes must be retained simultaneously?
  \end{description}

  \begin{block}{Notation}
    $b$ = branching factor, \qquad $d$ = depth of the shallowest goal, \qquad $m$ = maximum search depth.
  \end{block}
\end{frame}

\sectionframe{Part 3: Breadth-First Search}

\begin{frame}[shrink=20]{Breadth-First Search Expands by Depth}
  \begin{columns}[T,onlytextwidth]
    \begin{column}{0.52\textwidth}
      \begin{cpdefinition}
        \textbf{Breadth-First Search (BFS)} expands the shallowest unexpanded node first.
      \end{cpdefinition}

      \begin{itemize}
        \item Explore depth $0$, then depth $1$, then depth $2$, and so on.
        \item The frontier is maintained as a FIFO queue.
        \item First discovered $\rightarrow$ first expanded.
      \end{itemize}
    \end{column}

    \begin{column}{0.42\textwidth}
      % Search direction: BFS tree colored by depth layers, with queue ordering.
      \lectureimage[2.45in]{slide17.png}
    \end{column}
  \end{columns}
\end{frame}

\begin{frame}[shrink=20]{Breadth-First Search: General Structure}
  \begin{block}{Pseudocode}
    \ttfamily\small
    BFS(initial\_state, goal\_test):\\
    \quad frontier $\leftarrow$ FIFO queue containing initial\_state\\
    \quad explored $\leftarrow \emptyset$\\[0.04in]
    \quad while frontier is not empty:\\
    \qquad node $\leftarrow$ frontier.dequeue()\\
    \qquad if goal\_test(node.state): return solution(node)\\
    \qquad add node.state to explored\\
    \qquad for each action in ACTIONS(node.state):\\
    \qquad\quad child $\leftarrow$ RESULT(node.state, action)\\
    \qquad\quad if child is new: frontier.enqueue(child)\\
    \quad return failure
  \end{block}

  \begin{keyidea}
    FIFO ordering is what turns the generic search procedure into breadth-first exploration.
  \end{keyidea}
\end{frame}

\begin{frame}[shrink=20]{BFS: Guarantees and Cost}
  \begin{columns}[T,onlytextwidth]
    \begin{column}{0.48\textwidth}
      \begin{block}{Why minimum depth?}
        If the shallowest goal is at depth $d$, BFS finishes depths $0,1,\ldots,d-1$ before exploring anything deeper than $d$.
      \end{block}

      \begin{itemize}
        \item Complete: \textbf{Yes}, for finite branching factor.
        \item Minimum-depth solution: \textbf{Yes}.
      \end{itemize}
    \end{column}

    \begin{column}{0.48\textwidth}
      \begin{block}{Complexity}
        \begin{itemize}
          \item Time: $O(b^d)$
          \item Space: $O(b^d)$
        \end{itemize}
      \end{block}

      \begin{itemize}
        \item The frontier retains a broad layer of alternatives.
        \item Memory often becomes BFS's practical limiting resource.
      \end{itemize}
    \end{column}
  \end{columns}

  \begin{keyidea}
    BFS buys systematic minimum-depth exploration with exponential frontier growth.
  \end{keyidea}
\end{frame}

\sectionframe{Part 4: Depth-First and Iterative Deepening Search}

\begin{frame}[shrink=20]{Depth-First Search Commits to One Branch}
  \begin{columns}[T,onlytextwidth]
    \begin{column}{0.52\textwidth}
      \begin{cpdefinition}
        \textbf{Depth-First Search (DFS)} expands the deepest unexpanded node first.
      \end{cpdefinition}

      \begin{itemize}
        \item Follow one branch as deeply as possible.
        \item Backtrack when the branch cannot continue productively.
        \item The frontier is maintained as a LIFO stack or through recursion.
      \end{itemize}
    \end{column}

    \begin{column}{0.42\textwidth}
      % Search direction: DFS tree with one highlighted deep branch and a backtracking arrow.
      \lectureimage[2.45in]{slide21.png}
    \end{column}
  \end{columns}
\end{frame}

\begin{frame}[shrink=20]{Depth-First Search: General Structure}
  \begin{block}{Pseudocode}
    \ttfamily\small
    DFS(initial\_state, goal\_test):\\
    \quad frontier $\leftarrow$ LIFO stack containing initial\_state\\
    \quad explored $\leftarrow \emptyset$\\[0.04in]
    \quad while frontier is not empty:\\
    \qquad node $\leftarrow$ frontier.pop()\\
    \qquad if goal\_test(node.state): return solution(node)\\
    \qquad add node.state to explored\\
    \qquad for each action in ACTIONS(node.state):\\
    \qquad\quad child $\leftarrow$ RESULT(node.state, action)\\
    \qquad\quad if child is new: frontier.push(child)\\
    \quad return failure
  \end{block}

  \begin{keyidea}
    The generic structure is nearly identical to BFS; frontier discipline changes the behavior dramatically.
  \end{keyidea}
\end{frame}

\begin{frame}[shrink=20]{DFS Trades Guarantees for Memory}
  \begin{columns}[T,onlytextwidth]
    \begin{column}{0.47\textwidth}
      \begin{block}{Properties}
        \begin{itemize}
          \item Complete: \textbf{No}, in general.
          \item Minimum-depth solution: \textbf{No}.
          \item Time: $O(b^m)$.
          \item Space: approximately $O(m)$ in the simplified tree-search model.
        \end{itemize}
      \end{block}
    \end{column}

    \begin{column}{0.47\textwidth}
      \begin{block}{Failure modes}
        \begin{itemize}
          \item infinite or very deep branch;
          \item deep goal found before shallow goal;
          \item strong sensitivity to successor ordering.
        \end{itemize}
      \end{block}
    \end{column}
  \end{columns}

  \begin{keyidea}
    DFS keeps little active search state, but it can spend substantial effort in an unproductive region.
  \end{keyidea}
\end{frame}

\begin{frame}[shrink=20]{Depth Limits Lead to Iterative Deepening}
  \begin{columns}[T,onlytextwidth]
    \begin{column}{0.47\textwidth}
      \begin{block}{Depth-Limited Search}
        DFS with a maximum depth $\ell$.
        \begin{itemize}
          \item prevents uncontrolled descent;
          \item fails if every solution lies deeper than $\ell$.
        \end{itemize}
      \end{block}
    \end{column}

    \begin{column}{0.47\textwidth}
      \begin{block}{IDDFS}
        Repeat depth-limited DFS with:
        \[
          \ell = 0,1,2,3,\ldots
        \]
        until a goal is found.
      \end{block}
    \end{column}
  \end{columns}

  \begin{keyidea}
    Iterative deepening searches increasingly deep prefixes without committing BFS-level memory.
  \end{keyidea}
\end{frame}

\begin{frame}[shrink=20]{IDDFS Combines Useful Properties of BFS and DFS}
  \begin{columns}[T,onlytextwidth]
    \begin{column}{0.47\textwidth}
      \begin{block}{Properties}
        \begin{itemize}
          \item Complete: \textbf{Yes}.
          \item Minimum-depth solution: \textbf{Yes}.
          \item Time: $O(b^d)$.
          \item Space: $O(d)$ in the simplified presentation.
        \end{itemize}
      \end{block}
    \end{column}

    \begin{column}{0.47\textwidth}
      \begin{block}{Why repeated work is acceptable}
        Nodes near the root are expanded several times, but exponentially growing trees contain most nodes near the deepest explored level.
      \end{block}
    \end{column}
  \end{columns}

  \begin{keyidea}
    IDDFS accepts modest repeated computation to obtain BFS-like depth guarantees with DFS-like memory usage.
  \end{keyidea}
\end{frame}

\begin{frame}[shrink=20]{Three Frontier Policies}
  \begin{center}
    \renewcommand{\arraystretch}{1.35}
    \begin{tabular}{l|l|c|c|l}
      \textbf{Strategy} & \textbf{Expand next} & \textbf{Complete?} & \textbf{Min depth?} & \textbf{Tradeoff} \\\hline
      BFS & shallowest & Yes & Yes & high memory \\
      DFS & deepest & No$^*$ & No & low memory \\
      IDDFS & deepest within limit & Yes & Yes & repeated work
    \end{tabular}
  \end{center}

  \vspace{0.08in}
  \scriptsize $^*$DFS can be complete on appropriately finite graph-search formulations, but not in the general unbounded case used for comparison here.

  \begin{keyidea}
    Uninformed strategies differ not in what the goal means, but in how they organize systematic exploration.
  \end{keyidea}
\end{frame}

\sectionframe{Part 5: Is Uninformed Search Really Artificial Intelligence?}

\begin{frame}[shrink=20]{``But Isn't This Just an Algorithm?''}
  \begin{columns}[T,onlytextwidth]
    \begin{column}{0.48\textwidth}
      \begin{block}{Why it feels ordinary}
        \begin{itemize}
          \item BFS and DFS use queues and stacks.
          \item They are deterministic procedures.
          \item They do not learn from data.
          \item They are standard graph algorithms.
        \end{itemize}
      \end{block}
    \end{column}

    \begin{column}{0.48\textwidth}
      \begin{block}{The better question}
        What capability does the algorithm provide to an agent?
        \begin{itemize}
          \item represent alternatives;
          \item predict consequences;
          \item construct a goal-directed plan.
        \end{itemize}
      \end{block}
    \end{column}
  \end{columns}

  \begin{keyidea}
    Algorithms and artificial intelligence are not opposites; AI systems are implemented through algorithms.
  \end{keyidea}
\end{frame}

\begin{frame}[shrink=20]{Search Is a Form of Internal Simulation}
  \begin{columns}[T,onlytextwidth]
    \begin{column}{0.55\textwidth}
      \begin{block}{A search-based agent can ask}
        \begin{itemize}
          \item What happens if I take action $A$?
          \item What becomes possible afterward?
          \item What if I instead take action $B$?
          \item Does any sequence eventually reach the goal?
        \end{itemize}
      \end{block}

      \begin{itemize}
        \item The agent need not physically execute every alternative.
        \item The transition model predicts consequences inside the representation.
      \end{itemize}
    \end{column}

    \begin{column}{0.39\textwidth}
    \end{column}
  \end{columns}

  \begin{keyidea}
    Deliberation separates considering an action from executing it.
  \end{keyidea}
\end{frame}

\begin{frame}[shrink=20]{Why This Was Legitimately AI}
  \begin{itemize}
    \item Newell and Simon proposed that human problem solving could be modeled as symbolic manipulation in a problem space.
    \item Logic Theorist demonstrated that theorem proving could be implemented as search over partial proofs.
    \item The scientific hypothesis was not that ``graphs are intelligent.''
    \item It was that important forms of reasoning could emerge from representation, operators, goals, and systematic exploration.
  \end{itemize}

  \begin{block}{Agent-level interpretation}
    environment $\rightarrow$ representation $\rightarrow$ possible futures $\rightarrow$ search $\rightarrow$ plan $\rightarrow$ action
  \end{block}

  \begin{keyidea}
    Search belongs in AI because it operationalizes a theory of goal-directed reasoning.
  \end{keyidea}
\end{frame}

\begin{frame}[shrink=20]{Why Search Feels Less Like AI Today}
  \begin{itemize}
    \item Techniques often stop feeling like AI once they become well understood and widely standardized.
    \item Search is now taught in algorithms courses because its behavior can be characterized precisely.
    \item Familiarity does not erase the role it plays inside autonomous reasoning systems.
  \end{itemize}

  \begin{alertblock}{Important distinction}
    BFS by itself is not an intelligent agent. An autonomous agent that uses a state-space model and BFS to construct a goal-directed plan is performing a foundational form of AI problem solving.
  \end{alertblock}

  \begin{keyidea}
    The algorithm may be ordinary; the capability it enables---reasoning about alternatives before acting---is fundamental.
  \end{keyidea}
\end{frame}

\sectionframe{Part 6: Application --- Autonomous Micromouse}

\begin{frame}[shrink=20]{Application: The Micromouse Competition}
  \begin{columns}[T,onlytextwidth]
    \begin{column}{0.54\textwidth}
      \begin{itemize}
        \item Micromouse is an autonomous robotics competition centered on maze discovery and navigation.
        \item The robot carries its own computation, sensing, actuation, and power.
        \item It does not begin with a complete map of the competition maze.
        \item It must discover enough structure to reach the goal and exploit what it learns.
      \end{itemize}

      \begin{block}{AI question}
        How can an autonomous agent construct a useful route when its model of the environment is initially incomplete?
      \end{block}
    \end{column}

    \begin{column}{0.40\textwidth}
      % Search direction: top-down Micromouse competition photograph showing autonomous robot inside a maze.
      \lectureimage[2.48in]{slide33.jpg}
    \end{column}
  \end{columns}

  \note{[Sources] Veritasium, ``The Fastest Maze-Solving Competition On Earth,'' May 24, 2023, https\://www.youtube.com/watch?v=ZMQbHMgK2rw. The video describes autonomous maze-solving robots, local sensing, search algorithms, and later high-speed runs.}
\end{frame}

\begin{frame}[shrink=20]{Abstract the Robot into a Search Problem}
  \begin{columns}[T,onlytextwidth]
    \begin{column}{0.53\textwidth}
      \begin{description}
        \item[State] Current maze cell.
        \item[Initial state] Starting cell near a corner.
        \item[Goal] Reach a designated center goal cell.
        \item[Actions] Move north, south, east, or west when a passage is open.
        \item[Transition] Move to the adjacent cell in the selected direction.
      \end{description}
    \end{column}

    \begin{column}{0.41\textwidth}
      % Search direction: grid maze labeled S and G with N/S/E/W discrete actions.
      \lectureimage[2.42in]{slide34.jpeg}
    \end{column}
  \end{columns}

  \begin{keyidea}
    For this lecture, the physical robot becomes an agent moving among discrete states in a graph.
  \end{keyidea}
\end{frame}

\begin{frame}[shrink=20]{The Maze Is a Graph the Robot Must Discover}
  \begin{columns}[T,onlytextwidth]
    \begin{column}{0.48\textwidth}
      \begin{block}{Maze $\rightarrow$ graph}
        \begin{itemize}
          \item cell $\rightarrow$ state;
          \item open passage $\rightarrow$ transition;
          \item wall $\rightarrow$ unavailable action;
          \item center $\rightarrow$ goal.
        \end{itemize}
      \end{block}
    \end{column}

    \begin{column}{0.48\textwidth}
      \begin{block}{But the graph is initially incomplete}
        \begin{itemize}
          \item onboard sensors reveal nearby walls;
          \item the internal map is updated incrementally;
          \item newly discovered walls remove assumed transitions.
        \end{itemize}
      \end{block}
    \end{column}
  \end{columns}

  \begin{center}
    \large sense $\rightarrow$ update map $\rightarrow$ search $\rightarrow$ move $\rightarrow$ sense again
  \end{center}

  \begin{center}
    \lectureimage[0.95in]{slide35.jpeg}
  \end{center}

  \note{[Sources] Veritasium, ``The Fastest Maze-Solving Competition On Earth.'' The video emphasizes local sensing and incremental discovery of maze walls.}
\end{frame}

\begin{frame}[shrink=20]{Wall Following Is Primarily Reactive}
  \begin{columns}[T,onlytextwidth]
    \begin{column}{0.52\textwidth}
      \begin{itemize}
        \item A simple strategy is to keep one side of the robot adjacent to a wall.
        \item The next action depends mainly on immediately perceived wall geometry.
        \item This requires little explicit reasoning over alternative future routes.
        \item Certain maze structures can defeat wall-following as a general strategy for reaching the intended goal.
      \end{itemize}

      \begin{block}{Agent interpretation}
        percept $\rightarrow$ local rule $\rightarrow$ action
      \end{block}
    \end{column}

    \begin{column}{0.42\textwidth}
      % Search direction: wall-following robot tracing a perimeter or loop around the maze.
      \lectureimage[2.40in]{slide36.jpeg}
    \end{column}
  \end{columns}

  \begin{keyidea}
    Micromouse gives us a physical contrast between reflex behavior and model-based deliberation.
  \end{keyidea}

  \note{[Sources] Veritasium, ``The Fastest Maze-Solving Competition On Earth,'' discussion of early wall-following strategies and later maze designs.}
\end{frame}

\begin{frame}[shrink=20]{Micromouse as Depth-First Search}
  \begin{columns}[T,onlytextwidth]
    \begin{column}{0.55\textwidth}
      \begin{enumerate}
        \item At a junction, choose an unexplored branch.
        \item Follow it as deeply as possible.
        \item At a dead end, backtrack to the most recent junction with another option.
        \item Continue until a goal is found.
      \end{enumerate}

      \begin{alertblock}{Limitation}
        The first discovered goal may use many more moves than another route that DFS has not explored yet.
      \end{alertblock}
    \end{column}

    \begin{column}{0.39\textwidth}
    \end{column}
  \end{columns}

  \note{[Sources] Veritasium, ``The Fastest Maze-Solving Competition On Earth,'' comparison of depth-first and breadth-first maze exploration.}
\end{frame}

\begin{frame}[shrink=20]{Micromouse as Breadth-First Search}
  \begin{itemize}
    \item BFS explores reachable cells in order of increasing number of moves from the start.
    \item In our uniform-step abstraction, the first goal found is a minimum-move route.
    \item Once a maze graph is known, BFS is therefore a natural shortest-depth planner.
  \end{itemize}

  \begin{columns}[T,onlytextwidth]
    \begin{column}{0.47\textwidth}
      \begin{block}{Search computation}
        Expand states in memory and organize the frontier by depth.
      \end{block}
    \end{column}

    \begin{column}{0.47\textwidth}
      \begin{block}{Physical exploration}
        Driving down a corridor, returning, and exploring another branch requires real motion.
      \end{block}
    \end{column}
  \end{columns}

  \begin{center}
    \lectureimage[0.78in]{slide38.jpeg}
  \end{center}

  \begin{keyidea}
    The order in which a search algorithm expands states does not have to equal the physical trajectory of the robot.
  \end{keyidea}

  \note{[Sources] Veritasium, ``The Fastest Maze-Solving Competition On Earth,'' comparison of BFS with DFS and discussion of physical backtracking.}
\end{frame}

\begin{frame}[shrink=20]{Flood Fill: A Practical Micromouse Strategy}
  \begin{columns}[T,onlytextwidth]
    \begin{column}{0.52\textwidth}
      \begin{itemize}
        \item A common classical Micromouse strategy is flood fill.
        \item Goal cells receive distance value $0$.
        \item Increasing values are propagated outward through traversable neighbors.
        \item The mouse moves toward a neighboring cell with a smaller value.
      \end{itemize}

      \begin{block}{Interpretation}
        The labels encode graph distance to the goal under the robot's current map.
      \end{block}
    \end{column}

    \begin{column}{0.42\textwidth}
    \end{column}
  \end{columns}

  \note{[Sources] Veritasium, ``The Fastest Maze-Solving Competition On Earth,'' flood-fill explanation; University at Buffalo IEEE Micromouse documentation and CSULB Micromouse flood-fill documentation describe queue-based distance propagation over the maze grid.}
\end{frame}

\begin{frame}[shrink=20]{Flood Fill Revises the Plan as Walls Are Discovered}
  \begin{center}
    \Large
    optimistic map $\rightarrow$ encounter wall $\rightarrow$ record wall $\rightarrow$ recompute $\rightarrow$ continue
  \end{center}

  \begin{itemize}
    \item Unknown passages can initially be treated optimistically as traversable.
    \item A sensed wall invalidates an assumed transition in the internal graph.
    \item Distance labels are updated using the revised connectivity.
    \item The mouse continues following decreasing values toward the goal.
  \end{itemize}

  \begin{keyidea}
    The agent repeatedly combines perception, representation, and search rather than solving the problem only once.
  \end{keyidea}

  \note{[Sources] Veritasium video transcript summary: the mouse begins with an optimistic map, marks newly discovered walls, and updates its route; see also CSULB-EATS Micromouse flood-fill documentation.}
\end{frame}

\begin{frame}[shrink=20]{Micromouse Is More Than ``Run BFS on a Graph''}
  \begin{columns}[T,onlytextwidth]
    \begin{column}{0.54\textwidth}
      \begin{itemize}
        \item \textbf{Perception:} detect local walls.
        \item \textbf{Representation:} maintain an internal map.
        \item \textbf{Memory:} retain discovered structure.
        \item \textbf{Search:} reason over routes in the current map.
        \item \textbf{Planning:} select an intended movement sequence.
        \item \textbf{Action:} execute movement through the physical maze.
      \end{itemize}
    \end{column}

    \begin{column}{0.40\textwidth}
      % Search direction: closed-loop diagram sensors -> map -> search -> plan -> actuators -> maze.
      \lectureimage[2.32in]{slide41.jpeg}
    \end{column}
  \end{columns}

  \begin{keyidea}
    The search algorithm is one component; the intelligent system is the autonomous agent that uses representation and search to act.
  \end{keyidea}
\end{frame}

\begin{frame}[shrink=20]{What We Deliberately Ignore Today}
  \begin{columns}[T,onlytextwidth]
    \begin{column}{0.48\textwidth}
      \begin{itemize}
        \item velocity and acceleration;
        \item braking distance;
        \item wheel dynamics and slip;
        \item turning radius.
      \end{itemize}
    \end{column}

    \begin{column}{0.48\textwidth}
      \begin{itemize}
        \item localization and sensor error;
        \item different execution times for turns and straight segments;
        \item continuous or diagonal motion.
      \end{itemize}
    \end{column}
  \end{columns}

  \begin{block}{Today's abstraction}
    One legal move between adjacent cells $=$ one action.
  \end{block}

  \begin{keyidea}
    Today we ask which sequence of discrete states reaches the goal. The physical question of how quickly the robot can execute that route comes next.
  \end{keyidea}

  \note{[Sources] Veritasium's Micromouse video later discusses diagonals, high-speed turns, control systems, and routes whose physical execution time differs from their geometric length; those issues are intentionally deferred here.}
\end{frame}

\begin{frame}[shrink=20]{Uninformed Search: Complete Mental Model}
  \begin{columns}[T,onlytextwidth]
    \begin{column}{0.46\textwidth}
      \begin{block}{The agent knows}
        \begin{itemize}
          \item initial/current state;
          \item available actions;
          \item transition model;
          \item goal test.
        \end{itemize}
      \end{block}
    \end{column}

    \begin{column}{0.50\textwidth}
      \begin{block}{The agent does not know}
        \begin{itemize}
          \item which unexplored state is promising;
          \item which branch appears closer to the goal;
          \item which successor deserves domain-specific preference.
        \end{itemize}
      \end{block}
    \end{column}
  \end{columns}

  \begin{center}
    \large environment $\rightarrow$ abstraction $\rightarrow$ state space $\rightarrow$ search $\rightarrow$ plan $\rightarrow$ action
  \end{center}

  \begin{keyidea}
    Search is not AI because it is complicated; it is AI because it provides a computational mechanism for goal-directed deliberation.
  \end{keyidea}
\end{frame}

\begin{frame}[shrink=20]{Exit Ticket: Choose the Search Strategy}
  \begin{activity}
    For each situation, choose BFS, DFS, or IDDFS and justify your choice using completeness, minimum solution depth, and memory requirements.
    \begin{enumerate}
      \item The goal is known to be shallow and memory is plentiful.
      \item The state space may be extremely deep and memory is severely constrained.
      \item The shallowest goal is important, but its depth is unknown and BFS memory is unacceptable.
    \end{enumerate}
  \end{activity}

  \begin{keyidea}
    Outcome: you can formulate an uninformed search problem, explain how frontier policy changes exploration, and defend search as a foundational AI mechanism.
  \end{keyidea}
\end{frame}

\begin{frame}[shrink=20]{Next Time}
  \begin{itemize}
    \item The discrete search model tells us \emph{where} the agent should go.
    \item A physical robot must also determine \emph{how} to execute that route.
    \item Orientation, turning, motion constraints, sensing, and execution time can change which route is practically best.
  \end{itemize}

  \begin{keyidea}
    Next: move from abstract state transitions to physical motion and richer planning constraints.
  \end{keyidea}
\end{frame}

\end{document}
